In this manuscript we focus on the idea to combine
Newton's interpolation formula for forward differences~\eqref{prop:newton-series-for-monomial-reversed}
and the hockey-stick family identities (e.g~\eqref{prop:segmented-hockey-stick-identity}) to express the sums of powers seamlessly.

This particular idea is great, however it can be generalized even further.
Thus, the main idea is to utilize an interpolation formula for power $n^m$ in terms of \textit{abstract difference operator} $D(n^m)$
and binomial coefficients $\binom{f(n)}{k}$.

We are not limited to linear difference operators $\Delta, \nabla, \delta$, there is a complete theory
on dynamic differential equations on time-scales, founded by Bohner and Peterson~\cite{Bohner2001DynamicEO}.
For example, Jackson~\cite{jackson_1909} derivative is given by

\begin{align*}
    D_q(f(x)) = \frac{f(qx) - f(x)}{qx-x}
\end{align*}

Thus, finding interpolation formula for $f(x)=x^n$ in terms of $D_q(f(x))$ implies new formulas for sums of powers.

In general, let be an abstract interpolation formula for powers
\begin{align*}
    n^m = \sum_{k} \binom{f(n)}{k} D(n^m, k)
\end{align*}
Thus, the formula of sums of powers involves the abstract difference operator $D$ evaluated at some point $k$ and the hockey-stick
family identity over the binomial coefficients $\binom{n}{k}$
\begin{align*}
    \KnuthRFoldSum{1}{n}{m} = \sum_{k} D(n^m, k) \sum_{j \leq n} \binom{f(j)}{k}
\end{align*}
Similarly, for multifold sums of powers
\begin{align*}
    \KnuthRFoldSum{r}{n}{m} = \sum_{k} D(n^m, k) \binom{f(n+r)}{k}
\end{align*}
Many interpolation approaches involve rising factorials $x_{(n)}$, falling factorials $(x)_n$,
or regular factorials $n!$,
and thus can be expressed
in terms of binomial coefficients, because
\begin{align*}
    \frac{(x)_n}{n!} = \binom{x}{n}; \quad \frac{x_{(n)}}{n!} = \binom{x+n-1}{n}.
\end{align*}

In particular, Donald Knuth provides the formula for multifold sums of odd powers~\cite{knuth1993johann}
based on operator of central finite differences evaluated in zero,
and Newton's interpolation formula for central differences~\cite{kolosov_2025_18098442}
\begin{proposition}[Multifold sums of odd powers]
    \label{prop:knuth-sums-of-odd-powers}
    \begin{align*}
        \KnuthRFoldSum{r}{n}{2m-1}
        &= \sum_{k=1}^{m} \binom{n+k-1+r}{2k-1+r} \frac{1}{2k} \delta^{2k} 0^{2m} \\
        &= \sum_{k=1}^{m} \binom{n+k-1+r}{2k-1+r} (2k-1)! T(2m, 2k)
    \end{align*}
\end{proposition}
where $T(n,k)$ are the central factorial numbers of the second
kind, see~\cite[section 58]{steffensen1927interpolation} and~\cite[formula (10a)]{carlitz1963divided}
\begin{align*}
    T(n, k) = \frac{1}{k!} \delta^k 0^n = \frac{1}{k!} \sum_{j=0}^{k} (-1)^j \binom{k}{j} \left( \frac{k}{2} - j \right)^n
\end{align*}
Central factorial numbers of the second kind $T(n,k)$ were defined by Riordan
in~\cite[ch. 6.5, formula (24)]{riordan1968combinatorial}
by polynomial identity
\begin{lemma} [Riordan power identity]
    \begin{align*}
        n^m = \sum_{k=1}^{m} T(m,k) \centralFactorial{n}{k}
    \end{align*}
\end{lemma}
where $\centralFactorial{n}{k}$ are central factorials
$\centralFactorial{n}{k} = n \prod_{j=0}^{k-1} \left( n + \frac{k}{2} -j \right)$.
The sequence
\href{https://oeis.org/A008957}{\texttt{A008957}}
in the OEIS~\cite{sloane2003line} provides non-zero central factorial numbers of the second kind $T(2n,2k)$.

As future research direction,
the Knuth's formula~\eqref{prop:knuth-sums-of-odd-powers} utilizes the operator
of central finite differences evaluated in zero.
Thus, it is worth to investigate the existence of sums of odd powers involving the central differences
evaluated at an arbitrary integer point $t$,
similar to multifold sums of powers via Newton's series for central differences.
